\section{QRPsとサンプルサイズ設計}
\subsection{授業内容}
\subsubsection{科目の中でのこのコマの位置づけ}
ここでは帰無仮説検定の仮定を逸脱した場合に何が起こるか，悪しき研究実践の例を仮想的に体験することでその問題を理解する。
帰無仮説検定はデータに対して厳格なルールを決めてから実践するべきもので，いわばスポーツのゲームのような側面がある。
ラフプレーをするとゲームが成立しないように，科学におけるラフプレーは真実を見誤るという問題を引き起こす。とくにサンプルサイズや帰無仮説を事後的に変えることで，本来コントロールされているはずのものがそうならなくなってしまう。
これらについて，\textcite{kosugi2023}の第6章を参考に話を進める。
\subsubsection{コマ主題細目}
\begin{description}
	\item[心理学の再現性問題] 2000年ごろから指摘され始めた，心理学の再現性問題を紹介する。その一因として挙げられる心理統計法の誤用とQRPs(問題ある研究実践)の存在を知り，このことが科学としての心理学の根幹を揺るがす大問題であることを共有したい。資料として\textcite{Ikeda2016}を用いる。
	\item[QRPsの実践] QRPsを仮想空間上で実践して，どのような問題が生じるかを見ていく。サンプルサイズを徐々に増やしていく(「もう少し頑張ってデータを取ろう！」)ことがなぜ悪いのか，有意なところだけ報告する(帰無仮説の事後的な変更)などによって，タイプ1エラーの確率が制御できない様子を確認する。
	      \begin{flushright}
		      $\to$\textcite{kosugi2023}のP.6--8およびP.225--232.
	      \end{flushright}
	\item[非心分布をつかったサンプルサイズ設計] 以上のことから，帰無仮説検定を行う際は事前にすべての設定・ルールを定めておく必要があることがわかる。実践者が恣意的に決められるのはサンプルサイズであるから，サンプルサイズの設計が重要になってくる。ここで対立仮説が従う分布として非心分布を導入し，これに基づくサンプルサイズ設計の例をシミュレーションで確認する。
	      \begin{flushright}
		      $\to$\textcite{kosugi2023}のP.6--8およびP.238--224.
	      \end{flushright}
	\item[非心分布をつかわないサンプルサイズ設計] 二群の平均値差の検定のような，比較的簡単な実験計画においてはサンプルサイズ設計の理論も単純なほうである。とはいえ，非心分布など慣れない確率分布の利用に戸惑うこともあろう。そこで多少力技的ではあるが，乱数生成を繰り返すことで実際どの程度の差・効果がどの程度の割合(確率)で検出できるかを検証するシミュレーションを行う。この方法は非心分布がわからないときや，複雑な実験計画にであっても応用が可能である。
\end{description}
\subsubsection{キーワード}
\begin{itemize}
	\item 再現性問題
	\item QRPs
	\item 非心分布
	\item 例数設計
\end{itemize}
\subsection{授業情報}
\paragraph{コマの展開方法}
講義/遠隔/演習
\subsubsection{予習・復習課題}
\paragraph{予習}  一年時の「データ解析基礎」のテキストを復習しておくこと。第19講が本時に対応する。
\paragraph{復習} 分散分析の例数設計については，テキストを参考に自ら実践することが望ましい。